\documentclass[11pt,a4paper]{article}

% --- Encodage / langue ---
\usepackage[utf8]{inputenc}
\usepackage[T1]{fontenc}
\usepackage[french]{babel}
\usepackage{lmodern}

% --- Mise en page compacte ---
\usepackage[a4paper,margin=1.9cm]{geometry}
\usepackage{microtype}
\usepackage{enumitem}
\setlist{nosep,leftmargin=1.2em}
\usepackage{booktabs}
\usepackage{titlesec}
\titlespacing*{\section}{0pt}{7pt}{2pt}
\titleformat{\section}{\normalfont\large\bfseries}{\thesection.}{0.5em}{}
\setlength{\parindent}{0pt}
\setlength{\parskip}{2pt}

\usepackage{hyperref}
\hypersetup{hidelinks}

\title{\vspace{-1.2cm}\textbf{Livrable 3 --- Liste des tests et de leurs scores\\[2pt]
\large Analyse automatique de partitions musicales (OMR)}}

\author{\normalsize Thibault Bourdoiseau \& Léo Li --- SAE TIFO 1 / TIFO 2 / MLRF}
\date{\vspace{-8pt}\small \today}

\begin{document}
\maketitle
\vspace{-10pt}

Complément à la description et à l'évaluation du prototype (document \emph{main}). On liste les
\textbf{tests unitaires} (validation du code) puis les \textbf{scores} de la version retenue (V2) selon les
critères du livrable~2. Mesures sur CameraPrIMuS, split train/test disjoint~; PA/RA/NA en \%.

\section{Tests unitaires (\texttt{uv run pytest})}
$34$ tests automatiques valident chaque brique de la chaîne~: prétraitement (Otsu, médian, adaptatif),
portées (métriques, localisation des 5 lignes, suppression), segmentation (composantes connexes), détection
des têtes, descripteurs (Hu, projections), classifieurs de durée et de silence ($k$-NN), armure, export MIDI.
\textbf{Résultat~: $34/34$ tests passent.}

\section{Scores du prototype retenu (V2)}
Banc d'essai bout-en-bout ($n=800$, test~$=240$ incipits disjoints), durées par classifieur $k$-NN.

\begin{center}
\begin{tabular}{lccccc}
\toprule
Sous-ensemble & $N$ & PA stricte & PA diat. & RA & NA \\
\midrule
PrIMuS simple  & 61  & 83{,}7 & 85{,}7 & 77{,}4 & 74{,}3 \\
PrIMuS général & 240 & 66{,}1 & 83{,}7 & 71{,}4 & 56{,}2 \\
général + armure oracle & 120 & 82{,}1 & 84{,}3 & 69{,}4 & 66{,}2 \\
\bottomrule
\end{tabular}
\end{center}

PA stricte = hauteur avec altérations~; PA diat.\ = sans (l'écart mesure le coût de l'armure non lue). La
ligne \emph{oracle} borne le gain d'une armure parfaitement lue. Reproduction~:
\texttt{uv run python scripts/benchmark.py -{}-n 800} (et \texttt{-{}-key-oracle} pour l'oracle).

\medskip
\textit{Pour mémoire}, les scores de la baseline géométrique (V1, état de l'art §5) et le détail des essais
intermédiaires (retenus / abandonnés) sont dans \texttt{SUIVI.md}.

\end{document}
